\subsection{Report: Observe Arrivals and Script Execution}
\label{sec:report}

Where \code{watch} observes outgoing evaluator send requests,
\filepath{lib/report.abclr2} observes the receiving object's dispatch
cycle. It supplies a replacement metaobject generator,
\code{meta/report-gen}, that reports arrivals, dispatch attempts, and
completed script bodies. The public \code{report} macro copies the
application object using this generator and rebinds its name to the copy.
This is not an evaluator wrapper installed on the original object.

Read Section~\ref{sec:copy} first. The copy retains references to scripts
and the outer lexical environment, but has a new object identity and queue.
The original Copy library's state-initializer limitation also applies:
do not assume that \code{report} can instrument arbitrary stateful
objects on this image. Our first client deliberately has no local state.

\subsubsection{Start with the Public Macros}

The supplemental \filepath{tutorial/examples/13-report.abclr2} explicitly
loads Copy before Report, whose original source requires Copy. Its
application and saved alias are:

\wholelisting[firstline=5,lastline=11]{tutorial/examples/13-report.abclr2}

The \code{:silent} handler intentionally sends no application reply.
Start with a request that does reply:

\begin{codeblock}
\begin{lstlisting}
(inc "tutorial/examples/13-report" :use-source t)
(report report-subject)
(eq report-subject report-original)       ; => NIL
[report-subject <== [:echo 7]]            ; => 7, plus reports
(noreport report-subject)
(eq report-subject report-original)       ; => T
\end{lstlisting}
\end{codeblock}

The macro takes a symbol naming an object, not an arbitrary expression.
Its three keyword options, \code{:message}, \code{:begin}, and
\code{:end}, all default to \code{t}. For example,
\code{(report report-subject :message nil :end nil)} selects only begin
reports. Disabling all three options suppresses printed reports but
still creates a copy; it does not make \code{report} a no-op.

Here is the original macro's construction sequence:

\wholelisting[firstline=30,lastline=37]{lib/report.abclr2}

The copier creates a new object with \code{meta/report-gen}; the macro
retrieves its default reporter, wraps it in a filter, saves the pair
\code{(new-object . previous-object)} in \code{*reporting-objects*},
installs the filter, and finally rebinds the supplied name. Other aliases
still denote the previous object. Requests through \code{report-original}
therefore bypass this new reporting metaobject.

\subsubsection{Follow One Dispatch Cycle}

Compare the original receive and begin handlers:

\wholelisting[firstline=60,lastline=79]{lib/report.abclr2}

The reporting protocol is
\code{[:report metaObject type [Message Reply Sender]]}.
Its identity field is the application's \emph{metaobject};
\code{[den metaObject]} gives the application object. The final triple
retains the received payload, application reply destination, and sender.
These are not the evaluator's completion continuation.

For a matching request, follow three distinct boundaries:

\begin{enumerate}
  \item \code{:message}: the metaobject receives a \code{:message}
        envelope and reports it before enqueuing the triple. If dormant,
        it becomes active and sends itself \code{:begin}.
  \item \code{:begin}: it dequeues a triple and reports it \emph{before}
        looking for a matching script. Bindings from a successful match
        extend the object's state environment for body evaluation.
  \item \code{:end}: the evaluator calls the supplied \code{cont} after
        the body finishes. That continuation sends an end report and the
        metaobject's internal \code{:end} request. The latter either
        starts another queued request or makes the object dormant.
\end{enumerate}

Thus a begin report means a dispatch attempt, not proof that a script
matched. An end report means that this body returned through evaluator
completion, not that all downstream work finished. A script may reply
early and continue working; another may finish without replying at all.
In our \code{:silent} case, a past-type send produces all three events
although there is no application answer. Do not call it with \code{<==}
and expect the end report to satisfy the waiting caller.

Reports themselves are past-type messages. The event order for one
completed request is arrival, begin, end, but reports from different
objects may interleave. They add work and cannot establish the timing
or scheduling of an uninstrumented run.

\subsubsection{Record the Actual Protocol Fields}

The teaching file includes a non-reifying recorder. This keeps the
diagnostic handler simple and avoids recursively instrumenting it:

\wholelisting[firstline=13,lastline=20]{tutorial/examples/13-report.abclr2}

In a fresh session after loading the file, replace the copy's reporter:

\begin{codeblock}
\begin{lstlisting}
(report report-subject :message nil :begin nil :end nil)
[[meta report-subject] <= [:set-reporter report-recorder]]
[report-subject <== [:echo 7]]
[events := [report-recorder <== :events]]
(mapcar #'second events)                 ; => (:MESSAGE :BEGIN :END)
(mapcar #'(lambda (e) (first (third e))) events)
; => ((:ECHO 7) (:ECHO 7) (:ECHO 7))
\end{lstlisting}
\end{codeblock}

Each record is \code{(metaObject type triple)}. The lambda only transforms
its ordinary Lisp argument; it neither captures ABCL variables nor sends
messages. Pass context explicitly if a host helper needs it.

\code{:set-reporter} replaces the \emph{whole} reporter chain. The macro's
earlier keyword filter is no longer in that path, so the recorder receives
all three event types. To select only end events while retaining structured
records, install a new filter:

\begin{codeblock}
\begin{lstlisting}
[only-end := [standard-filter-gen <== [:new report-recorder
                         :message nil :begin nil :end t]]]
[[meta report-subject] <= [:set-reporter only-end]]
[report-recorder <== :reset]
[report-subject <== [:echo 8]]
(mapcar #'second [report-recorder <== :events])
; => (:END)
(noreport report-subject)
\end{lstlisting}
\end{codeblock}

The generator's message requires all three named fields in this order;
it is an object protocol, not a Common Lisp keyword-argument function.
The filter forwards enabled events unchanged and drops disabled ones.
It does not suppress application messages or change script execution.
Reporter installation uses \code{<=}: the setter promises no application
reply, so changing it to \code{<==} would introduce a wait with no answer.

\paragraph{Reading the default report.}
Prefer structured records when interpreting reply and sender identities.
The \code{default-reporter} receives \code{[M R S]} and passes those
fields to \code{default-reporting-function} in the same order.
Its \code{from} field prints the sender; \code{reply=} prints the
application reply destination, not the evaluation continuation.
Either field can be \code{nil}. Printed objects and messages are abbreviated by
\code{*print-length*} and \code{*print-level*}. Do not use the displayed
names or labels as a substitute for the recorded triple.

\subsubsection{Restoration and the Limits of Instrumentation}

\code{noreport} looks up the current object in \code{*reporting-objects*},
removes its association, and rebinds the supplied name to the saved object.
It does not destroy the reporting copy, remove other aliases, drain its
queue, or merge changes back into the previous object. A retained alias
to the copy still runs and reports after \code{noreport}.

Repeated \code{report} creates successive copies, not multiple reporting
layers around the same object. Each copy has its own reporting metaobject;
it does not forward dispatch through the previous copy's metaobject.
With an intact registry, successive \code{noreport} calls restore the
previous copy and then the original. This differs from Watch's repeated
installation limit (Section~\ref{sec:watch}). Do not confuse this name
history with a general facility for rolling back application effects.

For an active application, copying is not an atomic snapshot. Existing
references still target the old object. Pending queue entries preserve
their original reply and sender references; Copy installs them with
\code{:enqueue}, not \code{:message}, so they do not acquire new arrival
reports merely by being copied. They also do not start simply because
the copy was returned. Follow the explicit \code{:begin} rules in
Section~\ref{sec:copy}, and do not manually start an already active
metaobject. This book uses an idle, empty-state subject to keep those
separate issues out of the first experiment.

\paragraph{Check your understanding.}
Will \code{report-original} start reporting when the other name is rebound?
\emph{Answer:} No; it retains the old object identity. Does an end report
replace \code{!} for the application caller? \emph{Answer:} No; evaluator
completion and application replies have separate destinations. Does
\code{noreport} stop a reporting copy held under another alias?
\emph{Answer:} No; it restores one binding, not every reference.

\paragraph{What the checks establish.}
Fifty-nine bounded probes check identities, alias preservation, event
triples, all eight filter combinations, repeated copying/restoration,
completion without an application reply, and dispatch state.
The formatter checks cover arguments and output for all three event types
and \code{nil} reply/sender fields. Run \code{make -C tutorial report-check}; copies,
inputs and transcripts remain in \filepath{build/report-tests/}.
These are bounded observations, not certification of arbitrary stateful
copying or transparent instrumentation.
